\section{A two-shortcut accelerated realization}
\label{sec:acceleration}

The table \texttt{PA194AX11GDHF} applies the two guarded shortcuts from
the incoming AX11GH/GHS acceleration work to the new AX11GDH parent.
The technique is inherited; its application to this new command design
and the resulting table are new within the inspected corpus.
No source instruction or arithmetic command changes. The native edits are
\[
\begin{array}{c|c|c}
\text{row}&\text{parent}&\text{accelerated}\\\hline
(26,1)&(1,R,42)&(1,R,91)\\
(101,0)&(0,L,26)&(0,R,17).
\end{array}
\]
Both tables have 194 working states and 388 transitions.

\subsection{An invariant on arbitrary complete tapes}

Let $T(x)$ be the bit at position $x$, with head position $x$.
In every run beginning in native state 74, on any initial complete tape,
the following local assertions hold whenever their state conditions hold:
\begin{align*}
q=101&\ \Longrightarrow\ T(x-1)=0,\\
q=11&\ \Longrightarrow\ T(x-1)=0\ \text{or}\ T(x+1)=0,\\
q=26\ \text{and}\ T(x)=1
 &\ \Longrightarrow\ T(x+1)=T(x+2)=0.
\end{align*}
These assertions are initially vacuous. The only incoming row to 101 is
$(96,0)\mapsto(0,R,101)$. The incoming rows to 11 are
$(14,0),(39,0),(56,0)$, each writing zero and moving left, and
$(38,0)$, writing zero and moving right. Thus the appropriate neighbor
of the new head is zero. In the parent, the only incoming rows to 26 are
$(11,0)$ and $(101,0)$, both writing zero and moving left.
The latter can only leave 26 scanning zero, by the first assertion.
In the former, if the new scanned bit is one, the second assertion at
the old head forces the other neighbor to be zero. This gives the last
assertion. In the accelerated table, the incoming row from 101 is removed;
the same argument still applies. Inspection of these finite row lists
and quantifier-free induction on the run prove the invariant in iPRA.

\subsection{Exact three-step identities}

At state 26 reading one, the next two cells to the right are zero. The
parent takes
\[
 (26,x)\longrightarrow(42,x+1)\longrightarrow(13,x+2)
                  \longrightarrow(91,x+1),
\]
with no tape change. The child takes the same whole-tape endpoint in one
step. At state 101 reading zero, the cell immediately left is zero, so
the parent takes
\[
 (101,x)\longrightarrow(26,x-1)\longrightarrow(194,x)
                   \longrightarrow(17,x+1),
\]
again with no tape change. Native state 194 here is a working control;
the native halt is 254. The child takes the same endpoint in one step.
All intermediate controls are nonhalting.

The standalone \texttt{build\_fast.py} checks the exact parent-table hash,
all incoming rows, the six path transitions and all twelve possible
length-five windows satisfying the respective guards. The arbitrary
exterior is untouched. Other rows are copied literally. These finite
identities together with the invariant give a uniform proof, rather than
an inference from the blank prefix.

\subsection{A positive clock inside iPRA}

Write $P_s$ and $F_t$ for the parent and faster configurations, extending
halt by an absorbing configuration. Define $k(C)=3$ on either changed
guarded row and $k(C)=1$ otherwise, including halt. Put
\[
 \sigma(0)=0,\qquad
 \sigma(t+1)=\sigma(t)+k(F_t).
\]
Every function here is primitive recursive. Quantifier-free induction
using the exact identities proves
\[
 F_t=P_{\sigma(t)},\qquad t\le\sigma(t)\le3t.
\]
If $H_P(t)$ means that the parent has halted by clock $t$, and $H_F(t)$
has the analogous meaning, absorption therefore gives
\[
 \mathsf{iPRA}\vdash
 \bigl(H_P(t)\to H_F(t)\bigr)\ \wedge\
 \bigl(H_F(t)\to H_P(3t)\bigr).
\]
For the first implication, $\sigma(t)\ge t$ ensures that $P_{\sigma(t)}$
is already halted. For the second, the endpoint at $\sigma(t)\le3t$
is a halting witness. Existentially quantifying the clock gives relative
halting equivalence; intuitionistic propositional reasoning gives relative
nonhalting equivalence. This transfers the parent's stated arithmetic
calibration without any new PA proof compiler.

The complete blank startup reaches the same source frame in 57,870
steps, compared with 57,922. It uses the first shortcut twenty times and
the second six times, accounting for exactly 52 saved steps. The builder
compares complete configurations after every accelerated step and the
corresponding one or three parent steps throughout that prefix.
The separate independent physical replay also proves all 4,778
arbitrary-register source outcomes directly from this accelerated table.

